%===========================================================
%  CV — Amine BENALI
%  Compiler (depuis ce dossier) : pdflatex cv.tex   (deux fois)
%  ou : npm run cv  (voir package.json) — génère public/cv/Amine_Benali_CV.pdf
%===========================================================
\documentclass[10pt,a4paper]{article}

\usepackage[T1]{fontenc}
\usepackage[utf8]{inputenc}
\usepackage[french]{babel}
\usepackage{charter}                % police sobre et lisible (Bitstream Charter)
\usepackage[a4paper,margin=1.0cm,top=0.75cm,bottom=0.75cm]{geometry}
\usepackage{titlesec}
\usepackage{enumitem}
\usepackage{xcolor}
\usepackage{tabularx}
\usepackage{graphicx}
\usepackage[hidelinks]{hyperref}

% ---------- Couleurs ----------
\definecolor{ink}{HTML}{1A1D21}
\definecolor{accent}{HTML}{14365D}
\definecolor{muted}{HTML}{5B626B}
\color{ink}

% ---------- Mise en page ----------
\pagestyle{empty}
\linespread{0.96}
\setlength{\parindent}{0pt}
\setlength{\parskip}{0pt}
\frenchsetup{StandardLists=true}   % garde des puces classiques avec babel français

\titleformat{\section}
  {\large\bfseries\scshape\color{accent}}{}{0pt}{}[\vspace{1pt}{\color{accent}\titlerule[0.6pt]}]
\titlespacing*{\section}{0pt}{5pt}{2pt}

\setlist[itemize]{leftmargin=1.1em,itemsep=0pt,topsep=1pt,parsep=0pt,label=\textbullet}

% ---------- Commandes ----------
% \entree{Poste / diplôme}{Dates}{Organisation}{Lieu}
\newcommand{\entree}[4]{%
  \textbf{#1}\hfill{\small #2}\\
  {\itshape #3}\hfill{\small\itshape\color{muted} #4}\par}

% \projet{Nom}{Technologies}
\newcommand{\projet}[2]{%
  \textbf{#1}\hfill{\small\itshape\color{muted} #2}\par}

% Photo : \phototrue affiche photo.jpg (même dossier), \photofalse la retire
\newif\ifphoto
\phototrue

\begin{document}

%================ EN-TÊTE ================
\ifphoto
\begin{minipage}[c]{0.15\textwidth}
  \includegraphics[width=2.3cm]{photo.jpg}
\end{minipage}\hfill
\begin{minipage}[c]{0.84\textwidth}
\fi
{\fontsize{24}{26}\selectfont\bfseries AMINE BENALI}\\[3pt]
{\large\color{accent} Ingénieur d'État en IA \& Data Science \textbar{} Développeur Full-Stack}\\[4pt]
{\small
Alger, Algérie \textbar{} +213 697 41 26 20 \textbar{}
\href{mailto:amine.benali.dev@gmail.com}{amine.benali.dev@gmail.com}\\
\href{https://portfolio-amine-benali.vercel.app/fr}{portfolio-amine-benali.vercel.app} \textbar{}
\href{https://www.linkedin.com/in/amine-benali-4b0056256/}{linkedin.com/in/amine-benali-4b0056256} \textbar{}
\href{https://github.com/benamine19}{github.com/benamine19}}
\ifphoto
\end{minipage}
\fi

%================ PROFIL ================
\section{Profil}
Ingénieur d'État et titulaire d'un master en intelligence artificielle et science des données (ESTIN, Bac+5),
développeur full-stack spécialisé dans les applications métier. Depuis 2023, je conçois et livre pour des
clients réels des plateformes de réservation, d'e-commerce, de facturation, de gestion et de santé connectée.
Je souhaite approfondir l'architecture logicielle, la qualité et les méthodes de développement dans un Master en génie logiciel.

%================ FORMATION ================
\section{Formation}

\entree{Diplôme d'Ingénieur d'État et Master en Intelligence Artificielle \& Data Science (Bac+5)}{2021 -- 2026}%
{École supérieure en sciences et technologies de l'informatique et du numérique (ESTIN)}{Béjaïa, Algérie}
\begin{itemize}
  \item Projet de fin d'études d'ingénieur : \textbf{19/20} \quad\textbar\quad Mémoire de master : \textbf{17,5/20}
\end{itemize}


\entree{Baccalauréat Technique Mathématique, option génie électrique -- moyenne \textbf{17,71/20}}{2021}{Lycée Lakhdar Belmadani}{Sétif, Algérie}

%================ COMMUNICATION ================
\section{Communication scientifique}

\textbf{Engineering a Sovereign Healthcare 4.0 Ecosystem: A Microservices-based Framework for Real time IoMT Monitoring and Native French Prescriptive Analytics}\hfill{\small 30 mai 2026}\\
{\small K.~Amara, A.~Charef, \textbf{A.~Benali} (LITAN, ESTIN) -- communication orale, 4\textsuperscript{e} Séminaire national sur la technologie et l'industrie, Tlemcen.}\par

%================ EXPÉRIENCE ================
\section{Expérience professionnelle}

\entree{Fondateur \& Développeur Full-Stack}{2026 -- aujourd'hui}{Celerix Studio}{Alger, Algérie}
\begin{itemize}
  \item Sites web, applications métier et logiciels de gestion : cadrage, architecture, développement, déploiement.
\end{itemize}


\entree{Stagiaire Développeur Web}{Juillet -- septembre 2025}{Dotrea}{Alger, Algérie}
\begin{itemize}
  \item Conception et développement d'une application de gestion des employés : dossiers, tâches, absences (Next.js, PostgreSQL).
  \item Développement d'un tableau de bord de suivi du chiffre d'affaires et des indicateurs de performance (Python, BI).
\end{itemize}


\entree{Développeur Full-Stack Freelance}{2023 -- 2025}{Clients directs}{Algérie}
\begin{itemize}
  \item Plateformes e-commerce (EquiposDz, Smart Fix, CarPro DZ), de voyage (Evora Voyage) et B2B médicale (Mujmaa).
  \item Responsable de bout en bout : interfaces, API, authentification, bases de données, déploiement et maintenance.
\end{itemize}


\entree{Développeur Web}{2024 (6 mois)}{Micro Strategy Solutions DZ}{Algérie}
\begin{itemize}
  \item Développement d'applications web et d'API back-end, dont des solutions destinées à un ministère.
\end{itemize}


\entree{Développeur Web (engagement associatif)}{Sept. 2023 -- juin 2024}{Nexus, club scientifique de cybersécurité}{ESTIN, Béjaïa}
\begin{itemize}
  \item Développement d'outils internes ; projets sur le développement sécurisé et l'analyse de vulnérabilités web.
\end{itemize}

%================ PROJETS ================
\section{Projets sélectionnés}

\projet{Plateforme intelligente pour la santé connectée (PFE)}{Next.js, Express, WebRTC, PostgreSQL, Python, PyTorch}
\begin{itemize}
  \item \emph{Design of a Smart Connected Health Platform} (intitulé officiel) : téléconsultation vidéo sécurisée (WebRTC),
        dossier médical électronique et intégration d'objets connectés, sur une architecture microservices.
  \item Deep learning (risque clinique, anomalies, aide à la prescription). Développeur principal : architecture, back-end, IA.
  \item Réalisé en binôme dans le cadre du dispositif Diplôme-Startup ; travaux présentés en communication orale (ci-dessus).
\end{itemize}


\projet{ACMA -- système d'information de prêts et de couverture médicale}{React, Flutter, Express, Prisma, MySQL}
{\small Circuit de validation des prêts, prises en charge médicales, remboursements mensuels, portail employé et application mobile.}\par


\projet{Evora Voyage -- plateforme B2B / B2C d'agence de voyage}{Next.js, Express, Prisma, MySQL}
{\small Offres de voyage, Omra, billetterie aérienne, tableau de bord d'administration et suivi statistique.}\par


\projet{Plateforme SaaS de facturation}{Next.js, Node.js, PostgreSQL}
{\small Devis et factures multi-devises, gestion des clients et des paiements, contrôle d'accès par rôles.}\par


%================ COMPÉTENCES ================
\section{Compétences et langues}
\begin{tabularx}{\textwidth}{@{}>{\bfseries}p{3.6cm}X@{}}
Génie logiciel      & POO, UML, design patterns, microservices, tests unitaires/d'intégration, Scrum, revue de code \\
Front-end et mobile & React, Next.js, Material UI, Tailwind CSS, Flutter \\
Back-end et données & Node.js, Express, Prisma, Django, API REST, WebSocket, WebRTC, PostgreSQL, MySQL, MongoDB \\
IA et données       & Python, PyTorch, deep learning, analyse de données, business intelligence \\
Outils et langages  & JavaScript, TypeScript, Java, SQL, Git, GitHub, Docker, Linux, CI/CD, Postman \\
Langues             & Arabe (langue maternelle), français B2 (TCF, 477 points), anglais B2 \\
\end{tabularx}

\end{document}
